\documentclass[]{article}

\usepackage{amsmath}
\usepackage{multirow}
\usepackage{natbib}
\usepackage{graphicx} 


\begin{document}

Extracting the thermal Sunyaev-Zel'dovich (tSZ) effect, a crucial probe of galaxy cluster thermodynamics, from microwave sky maps is hampered by astrophysical foregrounds, most notably the spatially correlated Cosmic Infrared Background (CIB). We present a Multi-Frequency Wiener Filter (MWF) designed to optimally isolate the tSZ signal by incorporating the complete auto- and cross-frequency power spectra of all sky components, treating the CIB as a source of correlated noise rather than a signal to be deterministically nulled. Applying this framework to simulated Simons Observatory and Planck observations across six frequency channels from 90 to 857 GHz, we reconstruct the tSZ Compton-y map and evaluate its fidelity against a standard Internal Linear Combination (ILC) method using a matched-filter cluster detection pipeline. Our analysis demonstrates that by explicitly modeling the CIB's spatial correlations, the MWF effectively suppresses foreground-induced fluctuations that contaminate the ILC reconstruction, resulting in a cluster catalog with substantially higher purity. While the MWF introduces a predictable, scale-dependent suppression of the tSZ signal characteristic of an optimal linear filter, it yields a significantly tighter mass-observable relation with lower scatter. These findings highlight that leveraging the full statistical covariance of foregrounds is critical for robustly extracting faint cosmological signals and maximizing the scientific return from next-generation CMB surveys.

\end{document}
                